\chapter{核心流程与控制逻辑}

\section{身份鉴权与游客拦截门禁流程}
在 LingoBridge 系统的安全设计中，对于游客与注册用户界限的管理采用了前端高阶门禁拦截机制。在前端应用的整个生命周期中，鉴权上下文作为全局的状态中枢，在初始化阶段会自动尝试从本地加密持久化缓存中提取历史会话令牌。如果令牌存在且未过期，系统将发起验证请求并挂载用户角色；如果令牌缺失或校验失效，上下文会将当前用户角色置为游客身份。对于系统中所有要求特定操作权限的交互元素，开发团队设计了门禁守护高阶组件。该组件通过监听鉴权状态的实时变化来动态决策其包裹的子节点的行为。当处于游客状态的用户试图点击该组件包裹的按钮或输入框时，门禁守卫会自动截断点击事件的向上传递，在当前视图顶层强制拉起毛玻璃效果的登录浮窗，并锁定底层页面的滚动。直到游客成功输入凭证完成登录并使鉴权上下文的状态更新后，登录浮窗才会平滑淡出，同时门禁组件释放拦截，允许后续的交互正常进行。

\section{教师布置课时与作业的一对一映射生成流程}
为了降低系统的状态复杂性，并避免由于复杂的关联结构造成的数据冗余与一致性问题，系统在教师端创建直播和课时作业时，采用了强制性的一对一映射状态生成流程。当教师在课程管理视图中选定目标课程并点击创建直播会话时，前端会引导教师填写直播课时的标题、预期上课时间与配套的作业要求。在教师确认提交后，系统将启动级联状态生成逻辑。前端服务适配器将这一系列参数打包，向后端的级联创建接口发起单一的原子性调用。后端接收到请求后，开启事务锁以防止并发冲突，首先调用用户和会话查询组件确认教师资格，接着在内存中同时生成三个具备全局唯一主键的实体记录，包括直播课堂实例节点、该直播课堂唯一挂载的教学课时单元，以及与该课时单元深度绑定且完全对应的一对一口语作业节点。这三个实体在生成时，通过外键约束在底层形成了牢固的环状关联关系。一旦其中任何一个实体的创建由于文件写冲突而发生异常，后端将执行级联回滚，清除其余已临时生成的内存映射，从而在物理结构上杜绝了孤立直播间或无归属作业等脏数据的出现，也使得学生在后续进入直播间时能够无缝定位并提交配套作业。

\section{电子白板画笔涂鸦与多页画布控制逻辑}
直播电子白板的涂鸦同步与清理是直播课堂中对交互质量要求极高的控制逻辑。为了在多分辨率设备之间实现画笔痕迹的精准还原，电子白板摒弃了传统的基于绝对像素的坐标记录方式，而采用了全生命周期的归一化坐标转换逻辑。当用户在画布层上按下鼠标或触控笔并进行滑动时，白板监听器会捕获当前触控点相对于画布物理左上角的偏移量，随后除以当前画布容器的实际宽度与高度，计算出绝对介于零到一之间的浮点型相对坐标。这些相对坐标点序列与笔刷颜色、线宽信息一同被实时组装为墨迹事件包。在多页课件展示场景下，电子白板的笔迹数据以当前课件的活动页码作为唯一哈希键在内存状态树中进行维护。当教师端触发翻页信令时，白板控制器会根据目标页码，在内存状态树中切回对应的历史笔迹数组，并在下一帧完成画布上下文的重新绘制，而对于其他未激活页码的笔迹数据则予以在内存中完整保留。当教师点击画布清理按钮时，系统会触发局部画布重置机制，该机制将严格限制仅删除内存状态树中当前激活页码对应的笔迹数据，而非清除整堂课的所有笔墨，随即调用清除画布方法重置当前透明画布，从而保证了师生在非当前活动页上的历史绘制成果不会因局部清理而丢失。
